The modular representation of discrete event systems can be exploited by composition synthesis to mitigate the state explosion problem that monolithic synthesis suffers. In this paper we present how compositional synthesis can be applied to LTL. We implement compositional synthesis for GR(1) goals and show that with a very basic heuristic for selecting subplants the compositional synthesis already achieves significant gains with respect to a monolithic method, both in a native DES synthesis tool and using reactive synthesis tools via translations. 


Future work must look into avoiding intermediate state explosion by means of smarter heuristics that impact the order in which a plant LTS are composed. An implementation of the compositional framework for LTL requires defining a projection function (which can simply be an existential elimination of  literals) and a method for computing safe controllers. The latter can be trivialized by the conservative decision of not removing any states at all, to ensure that no winning states are lost. However, this is unlikely to provide good performance results. Thus, the challenge with full LTL is to find an effective way of removing a significant portion of subplant losing states.  

